\documentclass[11pt]{article}
\usepackage[margin=1in]{geometry}
\usepackage{amsmath,amssymb,amsthm}
\usepackage{algorithm}
\usepackage{algpseudocode}
\usepackage{booktabs}
\usepackage{hyperref}

\newtheorem{definition}{Definition}
\newtheorem{theorem}{Theorem}
\newtheorem{lemma}[theorem]{Lemma}
\newtheorem{property}{Property}

\title{\textbf{Mathematical Specification, Direction Preservation, and Norm Bounds of Gradient Clipping (ALGO-NN-28)}}
\author{Team Scaibu \\ \small Email: \texttt{jaydeep.9664920749@gmail.com} \\ \small Architecture and Core Engine Team}
\date{\today}

\begin{document}
\maketitle

\begin{abstract}
Recurrent neural networks, deep Transformers, and non-convex energy landscapes frequently encounter exploding gradients where step updates destabilize optimization. Global gradient norm clipping rescales the aggregate parameter gradient vector $\mathbf{g}$ by $\min\left(1, \frac{M}{\|\mathbf{g}\|_2}\right)$. We prove that norm clipping strictly preserves the gradient direction angle while establishing an upper bound on parameter displacement, provide full pseudocode matching the reference implementation, and derive complexity.
\end{abstract}

\section{Notation and Mathematical Preliminaries}
\label{sec:notation}

Let $\mathbf{g} \in \mathbb{R}^P$ denote the concatenated parameter gradient vector for $P$ neural network parameters. Let $M \in \mathbb{R}_{>0}$ denote the maximum allowed global $\ell_2$ norm threshold.

\subsection{Global Norm Rescaling Formulation}
\paragraph{Intuitive Meaning:} When gradients explode into enormous magnitudes, updating weights by learning rate $\eta \mathbf{g}$ causes numerical overflow and destroys learned features. Global norm clipping scales the entire gradient vector down so its total length does not exceed threshold $M$, while keeping its exact direction untouched.

\paragraph{Formal Mathematical Definition:}
Let $\|\mathbf{g}\|_2 = \sqrt{\sum_{i=1}^P g_i^2}$ denote the global Euclidean norm.
The clipped gradient $\tilde{\mathbf{g}} \in \mathbb{R}^P$ is defined as:
\begin{equation}
\label{eq:norm_clipping}
\tilde{\mathbf{g}} = \begin{cases} \mathbf{g} & \text{if } \|\mathbf{g}\|_2 \le M, \\ \frac{M}{\|\mathbf{g}\|_2} \mathbf{g} & \text{if } \|\mathbf{g}\|_2 > M. \end{cases}
\end{equation}
In elementwise value clipping with threshold $c > 0$:
\begin{equation}
\label{eq:value_clipping}
\tilde{g}_i = \max(-c, \min(c, g_i)).
\end{equation}

\paragraph{Worked Numerical Example:}
Let $\mathbf{g} = [3.0, 4.0]^\top$ and threshold $M = 2.5$.
Norm: $\|\mathbf{g}\|_2 = \sqrt{3.0^2 + 4.0^2} = \sqrt{9 + 16} = 5.0$.
Since $5.0 > 2.5$, scale factor $s = \frac{2.5}{5.0} = 0.5$.
Clipped vector: $\tilde{\mathbf{g}} = 0.5 \times [3.0, 4.0]^\top = [1.5, 2.0]^\top$.
New norm: $\|\tilde{\mathbf{g}}\|_2 = \sqrt{1.5^2 + 2.0^2} = \sqrt{2.25 + 4.0} = \sqrt{6.25} = 2.5 = M$.

\subsection{Summary of Symbols}
\begin{itemize}
    \item $P \in \mathbb{N}$: Number of parameter coordinates.
    \item $\mathbf{g} \in \mathbb{R}^P$: Unclipped gradient vector.
    \item $M \in \mathbb{R}_{>0}$: Global $\ell_2$ norm clipping threshold.
    \item $c \in \mathbb{R}_{>0}$: Elementwise coordinate clipping threshold.
    \item $\tilde{\mathbf{g}} \in \mathbb{R}^P$: Clipped output gradient vector.
    \item $\|\mathbf{g}\|_2 \in \mathbb{R}_{\ge 0}$: Global Euclidean norm.
\end{itemize}

\section{Problem Definition and Core Concepts}
\label{sec:problem_definition}

\subsection{Comparison: Norm vs. Value Clipping}
\begin{itemize}
    \item \textbf{Global Norm Clipping:} Preserves the exact gradient vector direction $\frac{\tilde{\mathbf{g}}}{\|\tilde{\mathbf{g}}\|_2} = \frac{\mathbf{g}}{\|\mathbf{g}\|_2}$. Preferred for deep networks.
    \item \textbf{Value Clipping:} Clamps each coordinate independently to $[-c, c]$, which changes the gradient vector direction when only some coordinates exceed the threshold.
\end{itemize}

\section{Algorithmic Specification}
\label{sec:algorithm}

Algorithm~\ref{alg:grad_clipping} specifies the evaluation routine matching \texttt{impl.py}.

\begin{algorithm}[H]
\caption{Gradient Norm and Value Clipping Routine}
\label{alg:grad_clipping}
\begin{algorithmic}[1]
\Require Gradient vector $\mathbf{g} \in \mathbb{R}^P$, maximum norm $M > 0$, optional value threshold $c$, mode $\mathrm{mode} \in \{\text{norm}, \text{value}\}$.
\Ensure Clipped gradients $\tilde{\mathbf{g}} \in \mathbb{R}^P$, original norm $N_{\mathrm{orig}}$, clipping indicator $\mathrm{was\_clipped}$.
\State $N_{\mathrm{orig}} \gets \sqrt{\sum_{i=1}^P g_i^2}$
\If{$\mathrm{mode} = \text{norm}$}
    \If{$N_{\mathrm{orig}} > M$}
        \State $scale \gets \frac{M}{N_{\mathrm{orig}} + 10^{-12}}$
        \State $\tilde{\mathbf{g}} \gets [g_i \cdot scale]_{i=1}^P, \quad \mathrm{was\_clipped} \gets \text{True}$
    \Else
        \State $\tilde{\mathbf{g}} \gets \text{list}(\mathbf{g}), \quad \mathrm{was\_clipped} \gets \text{False}$
    \EndIf
\Else
    \State $c_{\mathrm{val}} \gets c$ if $c \neq \text{None}$ else $M$
    \State $\tilde{\mathbf{g}} \gets [\max(-c_{\mathrm{val}}, \min(c_{\mathrm{val}}, g_i))]_{i=1}^P$
    \State $\mathrm{was\_clipped} \gets \exists i : |g_i| > c_{\mathrm{val}}$
\EndIf
\State \Return $\tilde{\mathbf{g}}, \quad N_{\mathrm{orig}}, \quad \mathrm{was\_clipped}$
\end{algorithmic}
\end{algorithm}

\section{Theoretical Guarantees and Direction Preservation}
\label{sec:guarantees}

\begin{theorem}[Strict Direction Preservation and Global Norm Invariant]
\label{thm:direction_preservation}
\textbf{Assumptions:} Let $\mathbf{g} \in \mathbb{R}^P \setminus \{\mathbf{0}\}$ and threshold $M > 0$.
\textbf{Guarantee:} The clipped gradient $\tilde{\mathbf{g}}$ satisfies $\|\tilde{\mathbf{g}}\|_2 \le M$ and has collinear unit direction $\frac{\tilde{\mathbf{g}}}{\|\tilde{\mathbf{g}}\|_2} = \frac{\mathbf{g}}{\|\mathbf{g}\|_2}$, with cosine similarity $\cos(\tilde{\mathbf{g}}, \mathbf{g}) = 1.0$.
\textbf{Proof:}
If $\|\mathbf{g}\|_2 \le M$, $\tilde{\mathbf{g}} = \mathbf{g} \implies \|\tilde{\mathbf{g}}\|_2 \le M$ and $\cos(\tilde{\mathbf{g}}, \mathbf{g}) = 1$.
If $\|\mathbf{g}\|_2 > M$, $\tilde{\mathbf{g}} = \frac{M}{\|\mathbf{g}\|_2} \mathbf{g}$.
Norm: $\|\tilde{\mathbf{g}}\|_2 = \left\|\frac{M}{\|\mathbf{g}\|_2} \mathbf{g}\right\|_2 = \frac{M}{\|\mathbf{g}\|_2} \|\mathbf{g}\|_2 = M \le M$.
Cosine angle:
\begin{equation}
\cos(\tilde{\mathbf{g}}, \mathbf{g}) = \frac{\tilde{\mathbf{g}}^\top \mathbf{g}}{\|\tilde{\mathbf{g}}\|_2 \|\mathbf{g}\|_2} = \frac{\frac{M}{\|\mathbf{g}\|_2} \mathbf{g}^\top \mathbf{g}}{M \|\mathbf{g}\|_2} = \frac{M \|\mathbf{g}\|_2^2}{M \|\mathbf{g}\|_2^2} = 1.0.
\end{equation}
\textbf{Limitations:} Norm clipping bounds step length but does not resolve underlying structural exploding dynamics caused by ill-conditioned recurrent weight matrices.
\end{theorem}

\section{Complexity Analysis}
\label{sec:complexity}

\begin{itemize}
    \item \textbf{Time Complexity:} $\mathcal{O}(P)$ operations to compute the sum of squares and rescale coordinates.
    \item \textbf{Space Complexity:} $\mathcal{O}(P)$ memory for the clipped gradient array.
\end{itemize}

\section{Worked Numerical Example}
\label{sec:worked_example}

Let $\mathbf{g} = [6.0, 8.0], M = 5.0$:
\begin{enumerate}
    \item Total norm: $\sqrt{6.0^2 + 8.0^2} = \sqrt{36 + 64} = \sqrt{100} = 10.0$.
    \item Since $10.0 > 5.0$, scale factor $= \frac{5.0}{10.0} = 0.5$.
    \item Clipped gradient: $[6.0 \times 0.5, 8.0 \times 0.5] = [3.0, 4.0]$.
    \item Indicator $\mathrm{was\_clipped} = \text{True}$.
\end{enumerate}

\section{Numerical Considerations and Edge Cases}
\label{sec:numerical}

\begin{itemize}
    \item \textbf{Zero-Division Protection:} Adding $\epsilon = 10^{-12}$ in the scale denominator prevents division by zero when $\|\mathbf{g}\|_2 = 0$.
\end{itemize}

\begin{thebibliography}{9}
\bibitem{pascanu2013} R.~Pascanu, T.~Mikolov, and Y.~Bengio, ``On the Difficulty of Training Recurrent Neural Networks,'' \emph{ICML}, 2013.
\end{thebibliography}

\end{document}
